% =============================================================================
% GOV-AXIOME (GOVERNANCE) - Vollständige Ausformulierung
% BCM 2.2 META-Modul
% =============================================================================

\section{Governance-Axiome (MA-GOV)}

Die GOV-Axiome definieren die Regeln, Constraints und ethischen Leitplanken.
Sie beantworten die Frage: \textit{Was ist erlaubt?}

%==============================================================================
% MA-GOV-01: Constraints existieren
%==============================================================================

\subsubsection*{MA-GOV-01: Constraints existieren}

\textbf{Statement:} Es gibt Beschränkungen für Verhalten.

\textbf{Formel:}
\begin{equation}
\Gamma \neq \emptyset
\end{equation}

\textbf{Beschreibung:}
Nicht alles ist erlaubt oder möglich. Constraints umfassen:
\begin{itemize}
    \item Physische Grenzen (Zeit, Raum, Körper)
    \item Rechtliche Grenzen (Gesetze, Verbote)
    \item Soziale Grenzen (Normen, Sanktionen)
    \item Ökonomische Grenzen (Budget, Ressourcen)
\end{itemize}

\textbf{BCM-Relevanz:}
Fundament des GOV-Moduls. Das BCM modelliert Verhalten unter Nebenbedingungen. REGULATE-Interventionen wirken durch Constraint-Manipulation.

\textbf{BEATRIX-Relevanz:}
BEATRIX muss Constraints kennen:
\begin{itemize}
    \item Welche Constraints limitieren die Persona?
    \item Sind Verhaltensoptionen feasible?
    \item Können Constraints gelockert werden?
\end{itemize}

\textbf{Konsequenz bei Fehlen:}
Alle Verhaltensweisen wären möglich. Keine Beschränkungen modellierbar. Regulierung hätte keine Grundlage.

\textbf{Validiert:} META

\textbf{Abhängigkeiten:} MA-ONT-06

\textbf{Literatur:} Becker (1968); North (1990)

%==============================================================================
% MA-GOV-02: Kontextabhängigkeit von Constraints
%==============================================================================

\subsubsection*{MA-GOV-02: Constraints sind kontextabhängig}

\textbf{Statement:} Erlaubtes hängt vom Kontext ab.

\textbf{Formel:}
\begin{equation}
\Gamma = f(\rho)
\end{equation}

\textbf{Beschreibung:}
Was in einem Kontext erlaubt ist, kann in einem anderen verboten sein:
\begin{itemize}
    \item Alkohol: erlaubt privat, verboten am Steuer
    \item Kleidung: angemessen im Büro ≠ am Strand
    \item Kommunikation: informell unter Freunden ≠ formell im Gericht
\end{itemize}

\textbf{BCM-Relevanz:}
Ermöglicht kontextspezifische Governance. Das 5-Ebenen-Modell (MA-ONT-26) strukturiert Kontexte hierarchisch.

\textbf{BEATRIX-Relevanz:}
BEATRIX prüft kontextspezifische Constraints:
\begin{itemize}
    \item Welche Regeln gelten in diesem spezifischen Kontext?
    \item Gibt es Kontextwechsel, die Constraints ändern?
    \item Kontextspezifische Interventionen designen
\end{itemize}

\textbf{Konsequenz bei Fehlen:}
Universelle Constraints. Kontextsensitivität von Regeln nicht abbildbar.

\textbf{Validiert:} META

\textbf{Abhängigkeiten:} MA-GOV-01, MA-ONT-04

\textbf{Literatur:} Ostrom (2005); Crawford \& Ostrom (1995)

%==============================================================================
% MA-GOV-03: Autonomie-Primat
%==============================================================================

\subsubsection*{MA-GOV-03: Autonomie-Primat}

\textbf{Statement:} Individuelle Autonomie hat Vorrang.

\textbf{Formel:}
\begin{equation}
\forall \Gamma: \min(\text{Einschränkung}) \quad \text{s.t. } \max(\text{Zielerreichung})
\end{equation}

\textbf{Beschreibung:}
Das BCM ist nicht paternalistisch. Autonomie ist ein Wert an sich:
\begin{itemize}
    \item Menschen haben das Recht auf eigene Entscheidungen
    \item Auch ``schlechte'' Entscheidungen sind zu respektieren
    \item Interventionen nur bei klarem Mehrwert und minimal invasiv
\end{itemize}

\textbf{BCM-Relevanz:}
Ethisches Leitprinzip. Begrenzt den Interventionsraum. Begründet die Präferenz für Nudges über Verbote.

\textbf{BEATRIX-Relevanz:}
BEATRIX respektiert Autonomie:
\begin{itemize}
    \item Keine manipulativen Interventionen empfehlen
    \item Transparenz über Interventionsziele
    \item Opt-Out-Möglichkeiten vorsehen
\end{itemize}

\textbf{Konsequenz bei Fehlen:}
Keine ethische Grenze für Interventionen. Paternalismus würde legitimiert.

\textbf{Validiert:} META

\textbf{Abhängigkeiten:} MA-GOV-01

\textbf{Literatur:} Mill (1859); Sunstein (2014)

%==============================================================================
% MA-GOV-04: Nicht-Normativität
%==============================================================================

\subsubsection*{MA-GOV-04: Nicht-Normativität}

\textbf{Statement:} Das BCM beschreibt, schreibt nicht vor.

\textbf{Formel:}
\begin{equation}
\text{BCM}: \text{is} \rightarrow \text{is}, \quad \text{nicht: is} \rightarrow \text{ought}
\end{equation}

\textbf{Beschreibung:}
Das BCM ist ein deskriptives, kein normatives System:
\begin{itemize}
    \item Es sagt, wie Menschen sich verhalten, nicht wie sie sich verhalten sollen
    \item Es ist wertfrei bezüglich der Inhalte
    \item Die Ziele der Intervention werden extern gesetzt
\end{itemize}

\textbf{BCM-Relevanz:}
Wichtige Abgrenzung. Das BCM kann für viele Ziele eingesetzt werden. Es ist ein Werkzeug, keine Weltanschauung.

\textbf{BEATRIX-Relevanz:}
BEATRIX ist agnostisch bezüglich Zielen:
\begin{itemize}
    \item Das Zielverhalten wird vom Nutzer vorgegeben
    \item BEATRIX optimiert für das vorgegebene Ziel
    \item Ethische Bewertung des Ziels ist Nutzer-Verantwortung
\end{itemize}

\textbf{Konsequenz bei Fehlen:}
Das BCM würde bestimmte Verhaltensweisen als ``richtig'' definieren. Ideologische Aufladung.

\textbf{Validiert:} META

\textbf{Abhängigkeiten:} -- (meta-ethisches Axiom)

\textbf{Literatur:} Weber (1917); Hume (1739)

%==============================================================================
% MA-GOV-05: Transparenz-Gebot
%==============================================================================

\subsubsection*{MA-GOV-05: Transparenz-Gebot}

\textbf{Statement:} Interventionen sollen transparent sein.

\textbf{Formel:}
\begin{equation}
\text{Transparency}(I) \rightarrow \max
\end{equation}

\textbf{Beschreibung:}
Menschen haben ein Recht zu wissen, dass und wie sie beeinflusst werden:
\begin{itemize}
    \item Offenlegung von Nudges
    \item Erklärung von Default-Setzungen
    \item Keine verdeckte Manipulation
\end{itemize}
Transparenz kann Wirkung reduzieren, ist aber ethisch geboten.

\textbf{BCM-Relevanz:}
Ethische Grenze für Interventionsdesign. ``Libertarian Paternalism'' verlangt Transparenz als Legitimationsbedingung.

\textbf{BEATRIX-Relevanz:}
BEATRIX empfiehlt transparente Interventionen:
\begin{itemize}
    \item Interventionen, die bei Offenlegung noch wirken, bevorzugen
    \item Warnung bei verdeckten Interventionen
    \item Transparenz-Score als Bewertungskriterium
\end{itemize}

\textbf{Konsequenz bei Fehlen:}
Verdeckte Manipulation wäre legitimiert. Vertrauensverlust bei Bekanntwerden.

\textbf{Validiert:} META

\textbf{Abhängigkeiten:} MA-GOV-03

\textbf{Literatur:} Thaler \& Sunstein (2008); Hansen \& Jespersen (2013)

%==============================================================================
% MA-GOV-06: Bereits in axiom_beispiele.tex
%==============================================================================

%==============================================================================
% MA-GOV-07: 8 Interventionstypen
%==============================================================================

\subsubsection*{MA-GOV-07: Acht Interventionstypen}

\textbf{Statement:} Es gibt acht grundlegende Interventionstypen.

\textbf{Formel:}
\begin{equation}
I \in \{\text{NUDGE, BOOST, THINK, INCENTIVE, REGULATE, DESIGN, SOCIAL, IDENTITY}\}
\end{equation}

\textbf{Beschreibung:}
Die acht Typen decken das Interventionsspektrum ab:
\begin{itemize}
    \item \textbf{NUDGE}: Entscheidungsarchitektur ändern
    \item \textbf{BOOST}: Kompetenz stärken
    \item \textbf{THINK}: Reflexion anregen
    \item \textbf{INCENTIVE}: Anreize setzen
    \item \textbf{REGULATE}: Regeln setzen/durchsetzen
    \item \textbf{DESIGN}: Physische Umgebung ändern
    \item \textbf{SOCIAL}: Soziale Normen aktivieren
    \item \textbf{IDENTITY}: Identität ansprechen
\end{itemize}

\textbf{BCM-Relevanz:}
Taxonomie für Interventionsplanung. Jeder Typ hat charakteristische Stärken und Grenzen.

\textbf{BEATRIX-Relevanz:}
BEATRIX klassifiziert und empfiehlt Interventionstypen:
\begin{itemize}
    \item Welcher Typ passt zur Persona und zum Kontext?
    \item Kombination verschiedener Typen planen
    \item Trade-offs zwischen Typen explizit machen
\end{itemize}

\textbf{Konsequenz bei Fehlen:}
Keine systematische Klassifikation von Interventionen. Ad-hoc-Interventionsdesign.

\textbf{Validiert:} META

\textbf{Abhängigkeiten:} MA-GOV-06

\textbf{Literatur:} Eigene Taxonomie, aufbauend auf Thaler \& Sunstein (2008); Hertwig \& Grüne-Yanoff (2017)

%==============================================================================
% MA-GOV-08: Freiheits-Hierarchie
%==============================================================================

\subsubsection*{MA-GOV-08: Freiheits-Hierarchie}

\textbf{Statement:} Interventionen haben unterschiedliche Freiheitsgrade.

\textbf{Formel:}
\begin{equation}
\text{Freedom}: \text{BOOST} > \text{THINK} > \text{NUDGE} > \text{DESIGN} > \text{SOCIAL} > \text{IDENTITY} > \text{INCENTIVE} > \text{REGULATE}
\end{equation}

\textbf{Beschreibung:}
Die Ranking ordnet Interventionen nach Autonomieeinschränkung:
\begin{itemize}
    \item BOOST: Erweitert Fähigkeiten, maximale Freiheit
    \item THINK: Aktiviert eigene Reflexion
    \item NUDGE: Sanfter Schubs, leicht überwindbar
    \item ...
    \item REGULATE: Verbot/Gebot, minimale Freiheit
\end{itemize}

\textbf{BCM-Relevanz:}
Operationalisierung des Subsidiaritätsprinzips (MA-GOV-06). Gibt klare Eskalationsreihenfolge vor.

\textbf{BEATRIX-Relevanz:}
BEATRIX verwendet die Hierarchie:
\begin{itemize}
    \item Empfehlungen mit Freedom-Score versehen
    \item Bei mehreren wirksamen Optionen: höhere Freiheit bevorzugen
    \item Trade-off Freiheit vs. Effektivität transparent machen
\end{itemize}

\textbf{Konsequenz bei Fehlen:}
Keine Ordnung von Interventionen nach Freiheit. Subsidiarität wäre nicht operationalisierbar.

\textbf{Validiert:} META

\textbf{Abhängigkeiten:} MA-GOV-07

\textbf{Literatur:} Eigenes Konstrukt; vgl. Sunstein (2016)

%==============================================================================
% MA-GOV-09: Eskalations-Prinzip
%==============================================================================

\subsubsection*{MA-GOV-09: Eskalations-Prinzip}

\textbf{Statement:} Von mild nach stark eskalieren.

\textbf{Formel:}
\begin{equation}
\text{Try}(I_n) \text{ before } \text{Try}(I_{n+1}) \quad \text{where } \text{Freedom}(I_n) > \text{Freedom}(I_{n+1})
\end{equation}

\textbf{Beschreibung:}
Die Interventionsleiter wird von oben nach unten abgearbeitet:
\begin{enumerate}
    \item Erst BOOST versuchen
    \item Wenn unwirksam: NUDGE
    \item Wenn unwirksam: INCENTIVE
    \item Nur als letztes: REGULATE
\end{enumerate}

\textbf{BCM-Relevanz:}
Prozedurale Norm für ethisches Interventionsdesign. Verhindert vorschnelle Regulierung.

\textbf{BEATRIX-Relevanz:}
BEATRIX empfiehlt Eskalationspfade:
\begin{itemize}
    \item Startpunkt: mildeste wirksame Intervention
    \item Fallback-Optionen bei Unwirksamkeit definieren
    \item Monitoring für Eskalationsentscheidung
\end{itemize}

\textbf{Konsequenz bei Fehlen:}
Interventionen könnten mit der härtesten Stufe beginnen. Unverhältnismäßigkeit.

\textbf{Validiert:} META

\textbf{Abhängigkeiten:} MA-GOV-08, MA-GOV-06

\textbf{Literatur:} Sunstein (2014)

%==============================================================================
% MA-GOV-10: 28 Interventions-Dimensionen
%==============================================================================

\subsubsection*{MA-GOV-10: 28 Interventions-Dimensionen}

\textbf{Statement:} Interventionen werden über 28 Dimensionen in 9 Kategorien charakterisiert.

\textbf{Formel:}
\begin{equation}
\text{Dim} = \{\text{Freiheit}, \text{Psychologie}, \text{Zeit}, \text{Reichweite}, \text{Kontext}, \text{Ethik}, \text{Implementation}, \text{Effekt}, \text{Dynamik}\}
\end{equation}

\textbf{Beschreibung:}
Jede Intervention kann auf 28 Dimensionen bewertet werden:
\begin{itemize}
    \item \textbf{Freiheit}: Restrictiveness, Opt-Out, Reversibility
    \item \textbf{Psychologie}: Target System (1/2), Mechanism, Awareness
    \item \textbf{Zeit}: Duration, Onset, Decay
    \item ...
\end{itemize}

\textbf{BCM-Relevanz:}
Ermöglicht präzise Charakterisierung und Vergleich von Interventionen. Systematische Planung statt Intuition.

\textbf{BEATRIX-Relevanz:}
BEATRIX verwendet die Dimensionen für:
\begin{itemize}
    \item Detaillierte Beschreibung empfohlener Interventionen
    \item Matching Intervention-Persona-Kontext
    \item Trade-off-Analyse zwischen Dimensionen
\end{itemize}

\textbf{Konsequenz bei Fehlen:}
Grobe Klassifikation nur nach Typ. Nuancen zwischen Interventionen nicht erfassbar.

\textbf{Validiert:} META

\textbf{Abhängigkeiten:} MA-GOV-07

\textbf{Literatur:} Eigenes Konstrukt

%==============================================================================
% MA-GOV-11: Distributive Gerechtigkeit
%==============================================================================

\subsubsection*{MA-GOV-11: Distributive Gerechtigkeit}

\textbf{Statement:} Kosten und Nutzen fair verteilen.

\textbf{Formel:}
\begin{equation}
\forall \sigma: \text{Burden}(\sigma) \propto \text{Capacity}(\sigma)
\end{equation}

\textbf{Beschreibung:}
Interventionen sollen nicht systematisch bestimmte Gruppen benachteiligen:
\begin{itemize}
    \item Wer profitiert, wer trägt die Kosten?
    \item Vulnerable Gruppen besonders schützen
    \item Progressivität: Stärkere Schultern tragen mehr
\end{itemize}

\textbf{BCM-Relevanz:}
Gerechtigkeitsprinzip für Interventionsdesign. Verhindert, dass Interventionen bestehende Ungleichheiten verstärken.

\textbf{BEATRIX-Relevanz:}
BEATRIX prüft distributive Effekte:
\begin{itemize}
    \item Wer ist von der Intervention betroffen?
    \item Sind vulnerable Gruppen überproportional belastet?
    \item Alternative Designs mit besserer Verteilung?
\end{itemize}

\textbf{Konsequenz bei Fehlen:}
Interventionen könnten systematisch ungerecht sein. Verstärkung von Ungleichheit.

\textbf{Validiert:} META

\textbf{Abhängigkeiten:} MA-GOV-03

\textbf{Literatur:} Rawls (1971); Sen (2009)

%==============================================================================
% MA-GOV-12: Reversibilität bevorzugen
%==============================================================================

\subsubsection*{MA-GOV-12: Reversibilität bevorzugen}

\textbf{Statement:} Reversible Interventionen sind vorzuziehen.

\textbf{Formel:}
\begin{equation}
\text{Prefer}(I_{\text{reversible}}) \text{ over } I_{\text{irreversible}}
\end{equation}

\textbf{Beschreibung:}
Interventionen sollten rückgängig machbar sein:
\begin{itemize}
    \item Wenn sie nicht wirken: Abbruch möglich
    \item Wenn unerwünschte Effekte: Korrektur möglich
    \item Experimentation wird ermöglicht
\end{itemize}
Irreversible Interventionen erfordern höhere Evidenzschwellen.

\textbf{BCM-Relevanz:}
Risikominimierung bei Unsicherheit. Ermöglicht iteratives Vorgehen (MA-GOV-20).

\textbf{BEATRIX-Relevanz:}
BEATRIX bewertet Reversibilität:
\begin{itemize}
    \item Reversibility-Score für jede Intervention
    \item Bei Unsicherheit: reversible Option wählen
    \item Warnung bei irreversiblen Empfehlungen
\end{itemize}

\textbf{Konsequenz bei Fehlen:}
Irreversible Interventionen ohne besondere Vorsicht. Keine Korrekturmöglichkeit bei Fehlern.

\textbf{Validiert:} META

\textbf{Abhängigkeiten:} MA-GOV-06

\textbf{Literatur:} Arrow \& Fisher (1974); Sunstein (2005)

%==============================================================================
% MA-GOV-13: Manipulation vermeiden
%==============================================================================

\subsubsection*{MA-GOV-13: Manipulation vermeiden}

\textbf{Statement:} Keine verdeckte, täuschende Beeinflussung.

\textbf{Formel:}
\begin{equation}
\text{If } \text{Hidden}(I) \land \text{Deceptive}(I) \rightarrow \text{Reject}(I)
\end{equation}

\textbf{Beschreibung:}
Manipulation = verdeckt + täuschend. Nicht jeder Nudge ist Manipulation:
\begin{itemize}
    \item Offener Nudge: OK
    \item Verdeckter, aber wahrhaftiger Nudge: Grauzone
    \item Verdeckter + täuschender Nudge: Manipulation = NICHT OK
\end{itemize}

\textbf{BCM-Relevanz:}
Ethische Grenze. Das BCM lehnt Manipulation ab, nicht Beeinflussung per se. Der Unterschied ist Transparenz und Wahrhaftigkeit.

\textbf{BEATRIX-Relevanz:}
BEATRIX prüft auf Manipulation:
\begin{itemize}
    \item Ist die Intervention bei Offenlegung noch vertretbar?
    \item Werden falsche Informationen verwendet?
    \item Manipulation-Flag bei Grenzfällen
\end{itemize}

\textbf{Konsequenz bei Fehlen:}
Keine Grenze zwischen Nudge und Manipulation. Ethische Kritik berechtigt.

\textbf{Validiert:} META

\textbf{Abhängigkeiten:} MA-GOV-05

\textbf{Literatur:} Sunstein (2016); Bovens (2009)

%==============================================================================
% MA-GOV-14: Bereits in axiom_beispiele.tex
%==============================================================================

%==============================================================================
% MA-GOV-15: Reaktanz beachten
%==============================================================================

\subsubsection*{MA-GOV-15: Reaktanz beachten}

\textbf{Statement:} Einschränkungen können psychologische Reaktanz auslösen.

\textbf{Formel:}
\begin{equation}
P(\text{Resistance}|\text{Restriction}) > 0
\end{equation}

\textbf{Beschreibung:}
Reaktanz ist die psychologische Gegenwehr gegen wahrgenommene Freiheitseinschränkung:
\begin{itemize}
    \item ``Verbotene Frucht'' wird attraktiver
    \item ``Trotzreaktion'': das Verbotene tun
    \item Besonders stark bei autonomieorientierten Personen
\end{itemize}

\textbf{BCM-Relevanz:}
Erklärt kontraproduktive Effekte von Regulierung. REGULATE kann Backlash erzeugen. Reaktanz ist besonders relevant bei IDN-starken Personen.

\textbf{BEATRIX-Relevanz:}
BEATRIX schätzt Reaktanz-Risiko:
\begin{itemize}
    \item Wie reaktanz-anfällig ist die Persona?
    \item Wie freiheitsbeschränkend ist die Intervention?
    \item Alternativen mit weniger Reaktanz-Risiko?
\end{itemize}

\textbf{Konsequenz bei Fehlen:}
Nur positive Interventionseffekte modelliert. Backlash-Risiko ignoriert.

\textbf{Validiert:} KNU

\textbf{Abhängigkeiten:} MA-GOV-01

\textbf{Literatur:} Brehm (1966); Steindl et al. (2015)

%==============================================================================
% MA-GOV-16: Ausweichverhalten beachten
%==============================================================================

\subsubsection*{MA-GOV-16: Ausweichverhalten beachten}

\textbf{Statement:} Regulierung kann zu Umgehung führen.

\textbf{Formel:}
\begin{equation}
P(\text{Evasion}|\text{REGULATE}) > 0
\end{equation}

\textbf{Beschreibung:}
Verbote werden nicht immer befolgt -- oft wird nach Schlupflöchern gesucht:
\begin{itemize}
    \item Steuerregeln → Steuervermeidung
    \item Geschwindigkeitsbegrenzung → Radarwarner
    \item Drogenverbot → illegaler Markt
\end{itemize}

\textbf{BCM-Relevanz:}
Begrenzt die Wirksamkeit von REGULATE. Bei hoher Evasion-Wahrscheinlichkeit sind andere Interventionstypen vorzuziehen.

\textbf{BEATRIX-Relevanz:}
BEATRIX schätzt Evasion-Wahrscheinlichkeit:
\begin{itemize}
    \item Wie einfach ist Umgehung?
    \item Wie hoch ist die Entdeckungswahrscheinlichkeit?
    \item Ist die Regulierung durchsetzbar?
\end{itemize}

\textbf{Konsequenz bei Fehlen:}
Perfekte Compliance angenommen. Regulierungs-Illusionen.

\textbf{Validiert:} WTX

\textbf{Abhängigkeiten:} MA-GOV-07

\textbf{Literatur:} Becker (1968); Slemrod (2007)

%==============================================================================
% MA-GOV-17: Opt-Out-Prinzip
%==============================================================================

\subsubsection*{MA-GOV-17: Opt-Out-Prinzip}

\textbf{Statement:} Austritt muss möglich sein (außer bei REGULATE).

\textbf{Formel:}
\begin{equation}
\forall I \neq \text{REGULATE}: \text{OptOut}(I) = \text{possible}
\end{equation}

\textbf{Beschreibung:}
Alle Interventionen außer REGULATE müssen Opt-Out erlauben:
\begin{itemize}
    \item Default kann überschrieben werden
    \item Nudge kann ignoriert werden
    \item Incentive kann abgelehnt werden
\end{itemize}
REGULATE ist die Ausnahme -- Gesetze gelten für alle.

\textbf{BCM-Relevanz:}
Operationalisiert Autonomie-Primat (MA-GOV-03). Unterscheidet ``libertarian paternalism'' von echtem Paternalismus.

\textbf{BEATRIX-Relevanz:}
BEATRIX stellt Opt-Out sicher:
\begin{itemize}
    \item Jede Intervention außer REGULATE hat Opt-Out-Weg
    \item Opt-Out muss praktisch möglich sein (nicht nur theoretisch)
    \item Warnung bei de-facto nicht-überwindbaren Defaults
\end{itemize}

\textbf{Konsequenz bei Fehlen:}
Alle Interventionen wären de facto Zwang. Verlust des libertären Elements.

\textbf{Validiert:} META

\textbf{Abhängigkeiten:} MA-GOV-03

\textbf{Literatur:} Thaler \& Sunstein (2008)

%==============================================================================
% MA-GOV-18: Informed Consent
%==============================================================================

\subsubsection*{MA-GOV-18: Informed Consent}

\textbf{Statement:} Informierte Zustimmung bei Interventionen.

\textbf{Formel:}
\begin{equation}
\text{Legitimate}(I) \rightarrow \text{Informed}(\sigma) \land \text{Consenting}(\sigma)
\end{equation}

\textbf{Beschreibung:}
Ideale Intervention: Die Zielperson versteht und akzeptiert sie:
\begin{itemize}
    \item Informed: Weiß, dass und wie sie beeinflusst wird
    \item Consenting: Stimmt der Beeinflussung zu
\end{itemize}
Bei öffentlichen Interventionen: implizite Zustimmung durch demokratische Legitimation.

\textbf{BCM-Relevanz:}
Höchster ethischer Standard. In der Praxis oft abgeschwächt, aber als Ideal wichtig.

\textbf{BEATRIX-Relevanz:}
BEATRIX bewertet Consent-Level:
\begin{itemize}
    \item Ist explizite Zustimmung möglich/nötig?
    \item Ist implizite Zustimmung ausreichend?
    \item Wie kann Informiertheit erhöht werden?
\end{itemize}

\textbf{Konsequenz bei Fehlen:}
Interventionen ohne jede Zustimmung legitim. Paternalismus pur.

\textbf{Validiert:} META

\textbf{Abhängigkeiten:} MA-GOV-05

\textbf{Literatur:} Faden \& Beauchamp (1986)

%==============================================================================
% MA-GOV-19: Interventionen habituieren
%==============================================================================

\subsubsection*{MA-GOV-19: Interventionen habituieren}

\textbf{Statement:} Interventionswirkung lässt über Zeit nach.

\textbf{Formel:}
\begin{equation}
\frac{\partial \text{Effect}(I)}{\partial \tau} < 0 \quad \text{für } \tau \rightarrow \infty
\end{equation}

\textbf{Beschreibung:}
Die meisten Interventionen verlieren Wirkung:
\begin{itemize}
    \item Nudges werden habituiert (``Blindheit'' für Hinweise)
    \item Incentives verlieren Neuheit
    \item Nur REGULATE und DESIGN bleiben stabil
\end{itemize}

\textbf{BCM-Relevanz:}
Erklärt Langzeit-Ineffektivität von Interventionen. Begründet Notwendigkeit von Iteration (MA-GOV-20).

\textbf{BEATRIX-Relevanz:}
BEATRIX modelliert Decay:
\begin{itemize}
    \item Erwartete Wirkdauer der Intervention
    \item Wann ist Refresh nötig?
    \item Decay-resistente Interventionstypen bevorzugen für Langzeiteffekte
\end{itemize}

\textbf{Konsequenz bei Fehlen:}
Konstante Interventionswirkung angenommen. Überschätzung von Langzeiteffekten.

\textbf{Validiert:} WTX

\textbf{Abhängigkeiten:} MA-ONT-20

\textbf{Literatur:} Allcott \& Rogers (2014)

%==============================================================================
% MA-GOV-20: Iteration notwendig
%==============================================================================

\subsubsection*{MA-GOV-20: Iteration notwendig}

\textbf{Statement:} Interventionen müssen iterativ angepasst werden.

\textbf{Formel:}
\begin{equation}
\text{Optimal}(I) \text{ requires Feedback-Loop}
\end{equation}

\textbf{Beschreibung:}
``Set and forget'' funktioniert nicht:
\begin{itemize}
    \item Erstimplementierung ist selten optimal
    \item Kontexte ändern sich
    \item Habituation erfordert Variation
\end{itemize}
PDCA-Zyklus: Plan-Do-Check-Act.

\textbf{BCM-Relevanz:}
Prozessprinzip. Interventionen sind keine einmaligen Events, sondern laufende Optimierung.

\textbf{BEATRIX-Relevanz:}
BEATRIX empfiehlt iteratives Design:
\begin{itemize}
    \item Monitoring-Metriken definieren
    \item Anpassungszeitpunkte festlegen
    \item A/B-Test-Optionen vorsehen
\end{itemize}

\textbf{Konsequenz bei Fehlen:}
Statische Interventionen. Keine Lernschleifen. Suboptimale Langzeitergebnisse.

\textbf{Validiert:} META

\textbf{Abhängigkeiten:} MA-GOV-19

\textbf{Literatur:} Deming (1986); Haynes et al. (2012)

%==============================================================================
% MA-GOV-21: Evaluation notwendig
%==============================================================================

\subsubsection*{MA-GOV-21: Evaluation notwendig}

\textbf{Statement:} Wirksamkeit muss gemessen werden.

\textbf{Formel:}
\begin{equation}
\forall I: \text{Measure}(\text{Effect}(I))
\end{equation}

\textbf{Beschreibung:}
Keine Intervention ohne Evaluation:
\begin{itemize}
    \item Wurde das Zielverhalten erreicht?
    \item War die Intervention ursächlich?
    \item Gibt es unbeabsichtigte Effekte?
\end{itemize}

\textbf{BCM-Relevanz:}
Wissenschaftliches Prinzip. Das BCM verlangt Evidenzbasierung, nicht nur Plausibilität.

\textbf{BEATRIX-Relevanz:}
BEATRIX integriert Evaluation:
\begin{itemize}
    \item KPIs für jede Intervention definieren
    \item Evaluationsdesign (RCT, Quasi-Experiment) vorschlagen
    \item Datenquellen identifizieren
\end{itemize}

\textbf{Konsequenz bei Fehlen:}
Intuitionsbasierte Interventionen. Keine Lernfähigkeit. Persistenz unwirksamer Maßnahmen.

\textbf{Validiert:} META

\textbf{Abhängigkeiten:} MA-GOV-20

\textbf{Literatur:} Campbell \& Stanley (1966); Shadish et al. (2002)

%==============================================================================
% MA-GOV-22: A/B-Testing bevorzugen
%==============================================================================

\subsubsection*{MA-GOV-22: A/B-Testing bevorzugen}

\textbf{Statement:} Experimentelle Evaluation wenn möglich.

\textbf{Formel:}
\begin{equation}
\text{RCT} > \text{Observational} \quad \text{für Kausalität}
\end{equation}

\textbf{Beschreibung:}
Randomisierte Kontrollversuche (RCT) liefern kausale Evidenz:
\begin{itemize}
    \item Zufällige Zuweisung eliminiert Selektionsbias
    \item Vergleich Treatment vs. Control
    \item Gold-Standard für Wirksamkeitsnachweis
\end{itemize}

\textbf{BCM-Relevanz:}
Methodisches Ideal. Wo RCTs nicht möglich: Quasi-experimentelle Designs, natürliche Experimente.

\textbf{BEATRIX-Relevanz:}
BEATRIX empfiehlt Testdesigns:
\begin{itemize}
    \item Ist ein RCT möglich?
    \item Sample-Size-Berechnung
    \item Alternative Designs bei RCT-Unmöglichkeit
\end{itemize}

\textbf{Konsequenz bei Fehlen:}
Korrelation = Kausalität-Fehler. Überschätzung von Interventionseffekten.

\textbf{Validiert:} META

\textbf{Abhängigkeiten:} MA-GOV-21

\textbf{Literatur:} Fisher (1935); Duflo et al. (2008)

%==============================================================================
% MA-GOV-23: Spillover messen
%==============================================================================

\subsubsection*{MA-GOV-23: Spillover messen}

\textbf{Statement:} Nicht-intendierte Effekte müssen erfasst werden.

\textbf{Formel:}
\begin{equation}
\text{Measure}(\text{Effect}_{\text{unintended}}) \text{ alongside } \text{Effect}_{\text{intended}}
\end{equation}

\textbf{Beschreibung:}
Interventionen haben oft unbeabsichtigte Nebenwirkungen:
\begin{itemize}
    \item Positiver Spillover: Andere Verhaltensweisen verbessern sich
    \item Negativer Spillover: Moral Licensing, Ausweichverhalten
    \item Spillover auf Nicht-Zielpersonen
\end{itemize}

\textbf{BCM-Relevanz:}
Ganzheitliche Evaluation. Eine Intervention, die das Zielverhalten verbessert, aber Schlimmeres verursacht, ist netto negativ.

\textbf{BEATRIX-Relevanz:}
BEATRIX prüft auf Spillover:
\begin{itemize}
    \item Welche anderen Verhaltensweisen könnten beeinflusst werden?
    \item Spillover-Metriken in Evaluation aufnehmen
    \item Warnung bei hohem negativen Spillover-Risiko
\end{itemize}

\textbf{Konsequenz bei Fehlen:}
Nur Zieleffekt betrachtet. Systemische Auswirkungen ignoriert.

\textbf{Validiert:} META

\textbf{Abhängigkeiten:} MA-ONT-22, MA-GOV-21

\textbf{Literatur:} Dolan \& Galizzi (2015)

%==============================================================================
% MA-GOV-24: Langfristeffekte beachten
%==============================================================================

\subsubsection*{MA-GOV-24: Langfristeffekte beachten}

\textbf{Statement:} Kurzfristige Wirkung $\neq$ Langfristige Wirkung.

\textbf{Formel:}
\begin{equation}
\text{Effect}(\tau = \text{short}) \neq \text{Effect}(\tau = \text{long})
\end{equation}

\textbf{Beschreibung:}
Interventionen können unterschiedliche Kurz- und Langzeiteffekte haben:
\begin{itemize}
    \item Initial stark, dann Decay (NUDGE)
    \item Initial schwach, dann Wachstum (BOOST)
    \item Langfristige Backlash-Risiken
\end{itemize}

\textbf{BCM-Relevanz:}
Temporale Differenzierung. Kurze Studien überschätzen möglicherweise langfristige Wirkung.

\textbf{BEATRIX-Relevanz:}
BEATRIX modelliert zeitliche Dynamik:
\begin{itemize}
    \item Erwartete Effektkurve über Zeit
    \item Langzeit-Follow-Up empfehlen
    \item Warnung bei bekannten Decay-Patterns
\end{itemize}

\textbf{Konsequenz bei Fehlen:}
Kurzfristeffekte extrapoliert. Systematische Überschätzung von Langzeitwirkung.

\textbf{Validiert:} META

\textbf{Abhängigkeiten:} MA-ONT-05, MA-GOV-21

\textbf{Literatur:} Allcott \& Rogers (2014); Brandon et al. (2017)

%==============================================================================
% MA-GOV-25: Governance evolviert
%==============================================================================

\subsubsection*{MA-GOV-25: Governance evolviert}

\textbf{Statement:} Regeln und Constraints passen sich an.

\textbf{Formel:}
\begin{equation}
\Gamma(\tau + \Delta\tau) = f(\Gamma(\tau), \text{Feedback})
\end{equation}

\textbf{Beschreibung:}
Governance ist nicht statisch:
\begin{itemize}
    \item Gesetze werden reformiert
    \item Normen wandeln sich
    \item Best Practices entwickeln sich
\end{itemize}
Das GOV-Modul selbst ist Gegenstand von Evolution.

\textbf{BCM-Relevanz:}
Meta-Dynamik. Die Regeln für Interventionen ändern sich. Ethische Standards entwickeln sich.

\textbf{BEATRIX-Relevanz:}
BEATRIX berücksichtigt Governance-Evolution:
\begin{itemize}
    \item Sind aktuelle Best Practices noch aktuell?
    \item Trends in der Interventionsforschung beobachten
    \item Empfehlungen können veralten
\end{itemize}

\textbf{Konsequenz bei Fehlen:}
Statische Governance-Annahmen. Keine Anpassung an veränderte Rahmenbedingungen.

\textbf{Validiert:} META

\textbf{Abhängigkeiten:} MA-ONT-25

\textbf{Literatur:} Ostrom (1990); North (1990)

\bigskip
\noindent\rule{\textwidth}{0.4pt}
\textit{Ende der GOV-Axiome. Alle 25 Governance-Axiome sind nun vollständig dokumentiert.}
